\chapter{Introduction}
\label{ch:introduction}

Many applications need the result of a query to stay up to date as the underlying database changes.
Recomputing that result from scratch after every change is wasteful, since each change usually affects only a small part of it.
\emph{Incremental view maintenance} (IVM) instead constructs a data structure that encodes the query result and maintains it incrementally as the database changes, so that the up-to-date result can be read from it when needed~\cite{olteanu2024IVM}.
How efficiently a query can be maintained this way depends on its structure and on the type of updates allowed.

\medskip

Abo~Khamis et al.~\cite{abokhamis2024Insert-Only} study the maintenance of \emph{full conjunctive queries}, join queries returning all attributes of the relations they combine.
We focus on the insert-only setting, in which tuples are only inserted and never deleted.
Such workloads arise in many practical settings, such as append-only event streams and continuously growing analytical datasets.
They show that in this setting a stream of inserts can be processed one tuple at a time in the same total time as evaluating the query once over the final database.
Maintaining the result incrementally is therefore essentially no more costly than computing it once.
Their approach organises this maintenance around a \emph{tree decomposition} of the query; we refer to it as the \emph{IVM$^+$ approach}.

Although the IVM$^+$ approach is theoretical in nature, in the insert-only setting it can be expressed directly as a set of SQL views.
The IVM$^+$ approach can therefore be realised on top of an existing IVM engine simply by rewriting the input query into these views, without modifying the engine itself.
This makes it lightweight to deploy, requiring only a query rewrite rather than the design or modification of a dedicated system.

\medskip

In this thesis, we investigate whether the IVM$^+$ approach is beneficial in practice when realised at the SQL level on an existing IVM engine.
We make two contributions:\footnote{The thesis repository is available at \url{https://github.com/ermin-mumic/UZH-Bachelor-Thesis-IVM}.}

\begin{itemize}
  \item \textbf{A query decomposer.} 
  We present a tool that, given a join query, computes a tree decomposition and produces the corresponding IVM$^+$ realisation as a set of SQL views.
  The entire translation is performed at the SQL level, without modifying the engine.
  The tool supports both the standard IVM$^+$ approach and an adjusted variant that we introduce in Chapter~\ref{ch:ivm}.
  It also supports aggregates such as \texttt{COUNT}, based on the concept of provenance semirings~\cite{green2007}, for both the standard approach and the variant.

  \item \textbf{An empirical evaluation on Feldera.}
  As our target engine we use Feldera~\cite{feldera2026}, a general-purpose engine based on the DBSP model~\cite{budiu2023DBSP} that incrementally maintains SQL views as the underlying data changes.
  In Feldera, queries are expressed as SQL views, which makes the rewritten IVM$^+$ views a natural fit.
  We compare the maintenance performance of the rewritten queries against that of the original queries.
\end{itemize}

Our evaluation finds that the rewrite helps most where the original query would otherwise build large intermediate results, on cyclic queries over dense data. There it maintains the result far more efficiently, in one case $42.96$ times faster than the original query and at the same memory, and it keeps tractable several queries the original cannot maintain within three hours. Where the original query is already cheap to maintain, the rewrite gives no speedup and uses up to $2.6$ times the memory, so its benefit is specific to these harder, cyclic queries rather than a property of the rewrite in general.

\medskip

The remainder of this thesis is organised as follows.
Chapter~\ref{ch:background} introduces the necessary background: the query model, the insert-only maintenance problem, tree decompositions and fractional hypertree width.
Chapter~\ref{ch:ivm} presents the IVM$^+$ approach, our adjusted variant, and a semiring extension for computing aggregates.
Chapter~\ref{ch:implementation} describes the query decomposer and the experiment harness, and how the approach is realised as Feldera views.
Chapter~\ref{ch:evaluation} reports the experimental setup, datasets, and results.
Chapter~\ref{ch:conclusion} concludes and discusses future work.
